% 本文件由 LaTeX 排版 Agent 根据有效 translation.json 自动生成。
% author=John Chrysostom; work_id=1903; title=致狄奥多尔的两次劝勉

\chapter{致狄奥多尔的两次劝勉}

% structure: p0001 source=h1 target=omitted reason=page-title-already-rendered-by-parent

% structure: p0002 source=h2 target=section reason=relative-heading-rank; base=h2

\section{第一封信}

\BibleQuote{但愿我的头变成水，我的眼成为泪泉！}\EditorialNote{耶 9:1} 如今我说这话正合时宜，是的，比当时的先知更合时宜。因为我虽不为多座城或整个民族哀悼，却要为与许多民族同等价值、甚至更为珍贵的灵魂哀悼。因为若一个遵行天主旨意的人比一万个违法者更好，那么你从前就比一万个犹太人更好。所以，若我编出比先知书更长的哀歌，发出更激烈的悲叹，无人能责备我。因为我哀悼的不是城市的倾覆，也不是恶人的被掳，而是一座神圣灵魂的荒凉，一座承载基督之殿的毁坏与抹除。因为凡曾在你荣耀之日认识你那被魔鬼点燃的井然心灵的人，岂不都会模仿先知的哀歌而悲叹？当他听见野蛮人的手玷污了至圣所，焚烧了一切：革鲁宾、约柜、赎罪盖、石版、金罐。这场灾祸比那更苦，是的，更苦，因为你灵魂中存放的抵押品远比那些更珍贵。这殿比那殿更神圣；因为它不是以金银发光，而是以圣神的恩宠；取代约柜和革鲁宾，它内住着基督、祂的圣父和\OriginalTerm{护慰者}{Paraclete}。但现在一切变了，殿已荒凉，失去昔日的美丽与俊秀，没有神圣而不可言喻的装饰，缺乏一切安全与保护；它没有门，没有闩，向各种毁灭灵魂可耻的思想敞开；若骄傲、邪淫、贪婪或比这些更可咒的思想想进入，无人阻挡；而从前，正如诸天对这一切不可进入，你灵魂的纯洁也是如此。如今我所说的，在那些现在目睹你荒凉与倾覆的人看来，也许难以置信；正因如此，我哀哭悲叹，且不停息，直到见你重振昔日光彩。因为这在人看来虽不可能，但为天主，一切都可能。因为祂\BibleQuote{从灰尘里抬举贫寒人，从粪堆中提拔穷乏人，使他们与王子同坐，就是与本国的王子同坐。}祂\BibleQuote{使不孕的妇女安居家中，成为多子的快乐母亲。}因此，不要对最完美的改变失望。因为若魔鬼有如此大能，将你从那美德的高峰巅顶坠入恶行的极端，那么天主岂不更能将你提回昔日的自信？不仅如此，甚至使你比从前更幸福。只是不要沮丧，不要抛弃美好的希望，不要陷入不虔敬之人的状态。因为使人陷入绝望的通常不是罪恶众多，而是灵魂的不敬。因此，撒罗满没有笼统地说：\Quote{\Emph{每一个}进入恶者窝巢的都藐视；}而只说：\Quote{不虔敬的人。}因为只有这类人进入恶者窝巢时才会如此。正是这种态度不让他们仰视，并重新攀登他们坠落的位置。因为这可咒的念头像轭压住灵魂的颈项，迫使它低头，阻止它仰望主人。勇敢优秀的人的本分就是打碎这轭，甩掉紧附身上的折磨者，说出先知的话：\BibleQuote{正如婢女的眼睛仰望主母的手，同样我们的眼睛仰望上主我们的天主，直到他怜悯我们。上主啊，求你怜悯我们，求你怜悯我们，因为我们已饱受轻蔑。}这些训诫真是神圣，是属于最高属灵智慧的旨意。经上说：我们饱受轻蔑，遭受了无数痛苦；但我们不会停止仰望天主，也不会停止向他祈祷，直到他接纳我们的恳求。因为高贵灵魂的标志就是不被压倒，不为重压的苦难所惊慌，也不多次祈祷无效而退却，而是坚持，直到他怜悯我们，正如圣达味所说。\par

2. 魔鬼之所以将我们投入绝望的思绪，是为了切断我们对天主的望德——这望德是安全的铁锚、我们生命的基石、引向天堂之途的向导、灭亡灵魂的救恩。经上说：\BibleQuote{我们因望德而得救}\BibleRef{罗 8:24}。的确，这望德如同从天上垂下的一根坚固绳索，支撑我们的灵魂，将紧紧抓住它的人逐渐引向至高之世，超拔他们脱离此生恶浪的暴风。任何人若变得软弱，松开这神圣的铁锚，便会立刻坠落，被淹没，陷入邪恶的深渊。那恶者深知此事，当他察觉到我们因恶行之知而自感压抑时，便亲自介入，将绝望带来的忧虑这重于铅块的重担加在我们身上；我们若接受此担，就必被重量立刻拖下，脱离那绳索，沉入苦难的深处——你自己如今正在那里，离弃了温良谦卑之主的诫命，执行那残忍暴君和救恩不共戴天之仇敌的一切命令；你打碎了容易的轭，抛开了轻省的重担，反以铁轭取而代之，是的，将沉重的磨石挂在颈上。当你淹没自己不幸福的灵魂，自陷于不断下沉的必然之中时，此后何处还能立足？那找到一块钱的妇人曾召邻人来分享喜乐，说：\Quote{与我同乐；}但我如今却要为我和你所有的朋友发出相反的邀请，不说\Quote{与我同乐}，而说\Quote{与我同哀}，和我同唱哀歌，同发悲鸣。因为最可怕的损失临到我身上，不是一定数量的金塔冷通，或大量宝石从我手中脱落，而是那位比这一切更宝贵、曾与我同航于这片浩瀚大海的人，不知如何竟失足落水，坠入了毁灭的深渊。\par

3. 若有人试图阻止我哀哭，我必以先知的话回答他们，说：\Quote{任凭我吧，我将痛哭；不要劳烦安慰我。} \BibleRef{依 22:4} 因为我现在的哀哭，并不是那种因过度哀恸而受谴责的；而是即使\Person{保禄}或\Person{伯多禄}在此，也不会羞于为此哭泣哀伤，并拒绝一切安慰的。因为凡为那人人共有的死亡而哀哭者，人们有理由指责他们心灵孱弱；但是，当一具尸体所替代的是一颗死去的灵魂，布满无数创伤，甚至在死亡中仍显露出它原先天生的俊美、健康和现已熄灭的美丽，谁能如此冷酷无情，以致用鼓励的话语代替哀号和痛哭呢？因为在另一个世界，不哀哭是神圣智慧的标志；同样，在这个世界，哀哭也是同一智慧的标志。那已升到高天、嘲笑现世虚幻、视身体之美不如石像、视黄金如泥土、视各种奢华如粪土的人，竟突然被狂热的欲望所吞没，毁坏了健康、刚强和青春的花朵，成了享乐的奴隶。那么，我请问你，我们岂不该为这样的人哭泣哀伤，直到我们把他带回吗？而这些事与人的灵魂有何相干呢？的确，在这个世界找不到从身体死亡中得释放的方法，但即便如此，也不能阻止哀恸者哭泣；然而，只有在这个世界，才可能使灵魂的死亡归于无效。\Quote{因为在阴府中}我们读到，\Quote{谁还会称谢你呢？} 那么，那些哀哭身体死亡的人，虽然知道他们的哀号不能使死者复活，却仍如此热切地哀哭，这不是极端的愚昧吗？而我们，明知常有希望将迷失的灵魂带回它从前的生命，却毫无表示，这岂不是更愚昧吗？因为许多人，无论是在现今还是在先祖的时代，虽已偏离正位，从正道中一头栽倒，却完全恢复了，以致后来的行为使先前的相形见绌，并获得了奖赏，戴上胜利的桂冠，被宣布为胜利者，列在圣人之中。因为只要一个人停留在享乐的火炉中，即使他眼前有无数这样的例子，他也觉得回改是不可能的；但一旦他开始从那里往外走一点点，继续前进，他就会把烈火最炽烈的部分抛在身后，看到前面的部分，他的脚步前满是露水和清凉。只是我们不要失望，不要厌倦回头；因为那受此影响的人，即使获得了无限的力量和热忱，也是徒然。因为他一旦把悔改之门向自己关闭，并堵住了进入赛场的入口，他如何能在外面成就任何好事，无论是小事还是大事呢？为此，\OriginalTerm{恶者}{Evil One}用各种诡计把（绝望的）念头种在我们里面；因为他若成功，就无需再费尽力气与我们搏斗；当我们匍匐在地、跌倒、不愿抵抗他时，他何须费劲呢？因为那能挣脱这条锁链的人，将恢复自己的力量，并会与魔鬼搏斗至最后一息；即使他摔倒无数次，他还会再站起来，打击他的敌人；但那被绝望的思虑所捆绑、使自己的力量松懈的人，他如何能得胜、如何能抵抗呢？他反倒逃跑了。\par

4. 且莫向我提及那些犯了小罪的人，但当设想一个充满一切邪恶的人，让他实践一切使他被排除在天国之外的事，并假设此人并非起初就不信者，而是一度属于信徒，且曾是悦乐天主的人，但后来成了奸淫者、通奸者、亲男色者、盗贼、酗酒者、鸡奸者、辱骂者，以及诸如此类的一切；我也不会赞同此人对自己绝望，即使他在如此巨大而不可言说的邪恶中一直活到极老。因为如果天主的义怒是一种\OriginalTerm{情欲}{passion}，那么人或许有理由绝望，无法扑灭他因作恶多端而点燃的火焰；但既然天主性体是\OriginalTerm{不动情}{passionless}的，即使祂惩罚，即使祂施以报复，祂这样做也不是出于忿怒，而是出于温柔的关怀和极大的仁爱；因此我们理当满怀勇气，信赖悔改的力量。因为即使那些得罪了祂的人，祂也不惯于为了自己而惩罚他们；因为没有任何伤害能穿透那天主性体；但祂行事是为我们的益处，并防止我们的乖戾因惯于藐视和忽略祂而变得更糟。因为正如一个人将自己置于光明之外，并未对光明造成损失，而是对自己造成最大的损失，被封闭在黑暗中；同样，一个惯于藐视那全能力量的人，并未对那力量造成伤害，而是对自己造成最大的伤害。为此，天主以惩罚来威胁我们，并时常施加惩罚，不是作为报复自己，而是为了吸引我们归向祂。因为一位医生也不会因那些神志不清者的侮辱而苦恼或恼怒，但仍然想方设法制止那些行为不当者，不谋自己的利益，而是为了他们的益处；如果他们表现出些许自制和清醒，他就欢喜快乐，更热切地施治，不是为以前的作为报复他们，而是希望增进他们的益处，使他们恢复完全健康的状态。同样，当我们陷入疯狂的极端时，天主言行一切，并非为报复我们以前的作为；而是因为祂希望将我们从失序中释放出来；而通过正确的理性，我们完全可以确信这一点。\par

5. 若有人就这些事与我们争辩，我们必以圣经来确认它们。试问，有谁比巴比伦王更为败坏？他虽曾深受天主大能的体验，以致向祂的先知俯首致敬，并命令向祂献上祭品和馨香，却又重陷于往昔的傲慢，将那些在天主面前不尊崇他的人捆绑投入火窑。然而，这个如此残暴不敬、近乎禽兽而非人者，天主仍召叫他悔改，赐予他多次回头机会：首先是火窑中发生的奇迹，其次是王所见而由达内尔解释的异象——一个足以软化铁石心肠的异象；除这些借由事件而来的劝诫外，先知也亲自劝告他说：\BibleQuote{因此，大王，请悦纳我的劝告：以施舍赎你的罪过，以怜悯穷人赎你的不义；或许你的过犯可得宽恕。}\BibleRef{达 4:27} 明智而蒙福的人，你作何言？如此大堕落之后，是否仍有归路？如此重病之后，是否还能痊愈？如此疯癫之后，是否仍有希望恢复清醒？这王已先自绝了一切希望：首先，他漠视了那位创造他、并引领他至此尊荣的天主，尽管他能列举无数发生在他自身及其祖先身上的大能与眷顾；然后，当他亲领天主智慧与预知的确切标记，并目睹魔术、天文及整套撒殚幻术体系被推翻后，却行出比先前更恶之事。因为智士、\OriginalTerm{加色丁人}{Gazarenes}所不能解释、且承认超越人性之事，一个被掳的青年却为他解明，并以此奇迹感动他，使他不仅自己相信，更成为普世这一教义的明晰宣导者与导师。\EditorialNote{参阅：达内尔书第二章} 因此，若未接受如此标记之前，他的无知已不可原谅，那么在此奇迹之后，在他宣认并将教导传播他人之后，更是无可宽恕。若他未曾真心相信祂是唯一真天主，他绝不会如此尊崇祂的仆人，也绝不会为他人制定如此律法。然而，在作出这种宣认之后，他再次堕入偶像崇拜：曾俯伏叩拜天主仆人的他，竟狂傲到将不崇拜他的天主仆人投入火窑。那么，天主是否按其所应得地惩罚了这个背教者？不！祂赐予他更伟大的大能标记，在他如此傲慢之后，又将他拉回昔日的景况；而更奇妙的是，为避免因奇迹过多而再度不信，祂行奇迹的对象正是王为捆绑投入火窑的孩童所点燃的火窑。仅仅熄灭火焰已是奇妙异常；但仁慈的天主为使他更生畏惧、增加惊惶、瓦解其铁石心肠，行了比这更大更奇的事：祂任凭火窑烧至王所愿的高度，然后彰显自己特有的能力——不是扑灭仇敌的计谋，而是在计谋实施时挫败它们。为免见他们幸免于火焰的人以为是幻象，祂让投他们入火窑的人被烧死，以此证明所见的确是火；否则它不会吞噬石脑油、麻屑、柴捆和如此众多的躯体。但没有任何事物比祂的命令更强；一切受造物的本性都服从那位从无中创造它们的主；这正是祂当时所彰显的：火焰虽容纳了可朽的躯体，却视之如不朽而远离它们，并安全地归还所受的托付，且更添光彩。犹如君王从宫廷出来，那些孩童也从火窑中走出，无人再耐心注视王，众目皆从王转向那奇景；王冠、紫袍及王室仪仗的一切华美，都不如那些忠信者在火中久留、然后如梦般走出者的景象更能吸引不信的群众。因为最脆弱的毛发——我们身体最脆弱的部分——竟比金刚石更持久地抵抗了吞噬一切的火焰。不仅他们被投入火中而毫发无损是奇迹，而且他们全程说话也是奇迹。凡见过火刑的人都知道，若紧闭双唇，尚可暂时抵挡烈焰；但若有人张口，灵魂即刻离躯。然而，如此大奇迹发生后，在场目睹者无不惊异，缺席者亦通过书信得知此事，那教导他人的王自己却毫无改善，重陷前恶。即使那时，天主仍未惩罚他，依然长久忍耐，借由异象和先知劝诫他。但当这一切都未能使他有所改善时，天主终于施加惩罚——并非为报复他先前所作，而是为断绝未来作恶的机会、遏止罪恶的进展；而且这惩罚并非永久性的，在责罚他数年之后，又恢复他昔日的尊荣，使他因惩罚未受任何损失，反而获得了最大的益处：对天主坚定的信德，以及为先前过犯的悔改。\EditorialNote{参阅：达内尔书第四章}\par

6. 天主的慈爱正是如此：祂从不拒绝真诚的悔改，即使有人堕入罪恶的深渊，却选择从那里回头，走向德行的道路，天主也接纳、欢迎，并竭尽全力使他恢复昔日的地位。祂所做的还更仁慈：即使有人没有表现出完全的悔改，祂也不忽略微小、不起眼的悔改，反而还为此赐予巨大的赏报；这从先知\Person{依撒意亚}论及犹太人民所说的话可以清楚看出，他这样说：\Quote{因他的罪，我使他痛苦片刻，击打他，转面不顾他；他痛苦哀伤，忧戚而行，然后我医治了他，安慰了他。}我们还可以举出另一个见证，即那个最不虔敬的君王，他由于妻子的影响而陷于罪恶：然而，当他只是悲伤、披上苦衣、并谴责自己的过犯时，他就如此为自己赢得了天主的仁慈，得以免除了那些迫临他的所有灾祸。因为天主对\Person{厄里亚}说：\Quote{你看见\Person{阿哈布}在我面前怎样心痛了吗？因为他在我面前哭泣，我就不在他的日子降灾祸。}此后再说到\Person{默纳舍}，他甚于所有人在狂怒与暴政，颠覆了合法的敬拜形式，关闭了圣殿，使偶像崇拜的欺骗盛行，变得比在他以前的所有人都更不虔敬；但后来他悔改了，就被列在天主的朋友中。如果他因看到自己罪孽深重而绝望于恢复与悔改，他就会失去后来所获得的一切：但事实上，他注视的是天主仁慈的无限，而非自己过犯的深重，并挣断了魔鬼的锁链，起来与它搏斗，完成了美好的赛程。\BibleRef{编下 33:10}天主不仅通过这些人的遭遇，也通过先知的话摧毁绝望的计谋，他这样说：\Quote{今天，如果你们听从祂的声音，不要再心硬，像在\Place{默黎巴}那样。}现在，\Quote{今天}这个词语可以在人生的任何时刻说出来，甚至老年临近时，只要你愿意：因为悔改不是以时间长短来衡量的，而是以灵魂的态度来衡量的。因为尼尼微人不需要很多天来涂抹他们的罪，而仅仅一天的时间就足以消除他们所有的过犯：那强盗也没有用很长时间来进入乐园，而是在可以说一句话的短暂瞬间，就洗净了他一生所犯的一切罪，并领受了天主的嘉奖，甚至在宗徒之前。我们也看到殉道者为自己获得荣冠，不是经过多年，而是经过数天，往往仅在一日之内。\par

7. 因此，我们在各方面都需要热忱，并且要有充分的心智准备：如果我们这样安排良心，憎恨从前的邪恶，并以天主所愿所命的热忱选择相反的道路，我们就不会因时间短暂而有所欠缺；至少许多后来者远超先到者。因为跌倒并不可怕，可怕的是跌倒后仍匍匐在地，不肯起来；并且怯懦懒惰，以绝望的推论掩盖道德意志的软弱。为此，先知也困惑地说：\BibleQuote{跌倒的岂能不起来？转离的岂能不回返？} \BibleRef{耶 8:4}但如果你问我关于信德后跌倒之人的实例，这一切话都是针对这样的人说的，因为跌倒的人原属于站立的人，不属于匍匐的人；否则处于那种状态的人怎能跌倒？但其他事也将藉比喻和更清楚的行为与言语说明。那隻从九十九隻中分离的羊\BibleRef{路 15:4}后来被领回，无非是向我们表明信友的跌倒与回头；因为它不是某个异类羊群的羊，而是属于同一数目，曾由同一位牧人牧放，它并非寻常迷路，而是流连于山岭和山谷，即一段远离正道的长途。难道牧人任凭它迷路吗？绝不，而是领它回来，既不驱逐，也不鞭打，却将它扛在肩上。因为正如最好的医生以精心的治疗使垂危的病人恢复健康，不仅按医术法则治疗，有时也给予满足：同样，天主引导那些极其堕落的人迈向德行，不是以严厉，而是柔和渐进，并从各方面扶持，使分离不致扩大，错误不致延长。同样的真理也蕴含在荡子的比喻中。因为他也不是外人，而是儿子，是那位讨父亲喜悦的孩子的兄弟，他陷入的并非寻常恶习，而是到了罪恶的极处，这富有、自由、有教养的儿子竟沦落到比家仆、外人、雇工更悲惨的境地。然而他回到原来的状况，恢复了从前的尊荣。但如果他绝望了生命，因所遭遇的而沮丧，留在异乡，他就不会得到所得到的，而会饿死，遭受最可怜的死亡。但他悔改了，没有绝望，甚至如此败落之后，仍恢复了从前的光彩，穿上最美的袍子，享受比未跌倒的兄长更大的荣耀。因为他说：\BibleQuote{这些年来，我服事你，从未违背过你的命令，而你从未给我一只山羊羔，让我和朋友一同欢乐；但你这个儿子，把家产和娼妓挥霍一空，却为他宰了肥牛犊。} \BibleRef{路 15:29}悔改的力量如此伟大。\par

8. 既然有如此伟大的榜样，让我们不要继续作恶，也不要对和好绝望，而是让我们自己也说：\Quote{我要到我父亲那里去}，并让我们亲近天主。因为祂从不远离我们，而是我们远离了祂：因为我们读到：\Quote{我是近处的天主，不是远处的天主。}再者，当祂藉这位先知的口责备他们时，祂说：\Quote{岂不是你们的罪孽使你们与我分离吗？}因此，既然这是使我们远离天主的原因，让我们移除这可憎的障碍，它阻止任何亲近的接近。\par

但是，现在请听这在真实事例中是如何发生的。在\Place{格林多}人中，有一位显赫的人物犯了一个罪，这罪甚至在外邦人中都没有被提及。此人是信徒，属于基督的大家庭；有人说他实际上还是司铎团的一员。那么，\Person{保禄}做了什么？他把他从那些走向救恩的人的共融中切断了吗？绝不：因为正是\Person{保禄}本人再三斥责\Place{格林多}人，因为他们没有带领那人悔改；但是，为了向我们证明没有不能治愈的罪，他再次论及那比外邦人犯罪更严重的人说：\BibleQuote{将这样的人交给撒殚，败坏他的肉体，好使他的灵魂在主耶稣基督的日子得以得救。}\BibleRef{格前 5:5}这是在悔改之前；但在他悔改之后，他说：\BibleQuote{这样的人，受了大多数人的这惩罚，已经够了。}\BibleRef{格后 2:6}并且他通过一封信吩咐他们再次安慰那人，接纳他的悔改，免得他被撒殚胜过。此外，当整个\Place{迦拉达}民族在信德之后跌倒，且为了对基督的信德曾行过奇迹、忍受过许多考验时，\Person{保禄}又使他们站起来。因为他们曾行过奇迹，他作证说：\BibleQuote{那赐给你们圣神，并在你们中间行奇迹的，是他吗？}\BibleRef{迦 3:5}而且他们为了信德忍受了许多争斗，他也作证说：\BibleQuote{你们受了许多苦，如果真是徒然，难道果真白费了吗？}\BibleRef{迦 3:4}然而，在取得如此大的进步之后，他们犯了足以使他们与基督隔绝的罪，关于这罪他宣告说：\BibleQuote{看，我\Person{保禄}告诉你们，如果你们受割损，基督就对你们毫无益处了。}又说：\BibleQuote{你们这些靠法律成义的人，是从恩宠上跌落了。}然而，即使在如此大的跌倒之后，他仍接纳他们，说：\BibleQuote{我的小孩子们，我为你们再受产痛，直到基督在你们内形成为止。}\BibleRef{迦 4:19}这表明，在极端堕落之后，基督可能再次在我们内形成：因为他不愿罪人死亡，而宁愿他悔改并生活。\par

9. 那么，让我们归向祂，我的挚友，并承行天主的旨意。因为祂创造了我们，使我们存在，为使我们分享永恒的福乐，为将天国赐予我们，而非将我们投入地狱、交付给烈火；因为这火不是为我们准备的，而是为魔鬼；但天国从古时已为我们预定和预备妥当。为表明这两项真理，祂对右边的人说：\Quote{我父所祝福的，你们来承受自创世以来为你们预备的国度吧！}但对左边的人说：\Quote{可咒骂的，离开我，到那永火里去吧！}（他不再说\InnerQuote{为你们}，而是说）\Quote{为魔鬼和他的使者预备的。}\BibleRef{玛 25:34}可见地狱不是为我们，而是为他及其使者预备的；但天国在创世之前已为我们预备好了。那么，我们不要使自己不配进入婚筵：因为只要我们还在这个世界，即使犯了无数罪恶，也能藉着为过犯表现悔改而全部洗净；但一旦我们离世到另一个世界，即使我们表现出最恳切的悔改，也无济于事，即使我们咬牙切齿、捶胸顿足、发出无数求救的呼声，也没有人会用指尖滴一滴水到我们灼烧的身体上，我们只会听到那富翁在比喻中所说的那句话：\Quote{在我们与你们之间，隔着一个巨大的深渊。}\BibleRef{路 16:26}那么，我恳求你，在此地清醒过来，让我们按应有的方式认识我们的主。因为只有当我们到了阴府，才应放弃从悔改而来的希望；在那里，这药方是软弱无用的；但在此地，即使在老年时才使用它，也显示出极大的力量。为此，魔鬼千方百计要在我们心中扎根绝望的推论：因为他知道，只要我们稍作悔改，便不会毫无报偿。正如那给一杯凉水的人，他的赏报已被保留；同样，那为自己所作的恶而悔改的人，即使他无法表现出与其过犯相称的悔改，也将获得相称的赏报。因为公义的审判者不会忽略任何善行，无论多么微小。如果祂如此严格地审查我们的罪，连我们的言语和思想都要惩罚，那么我们的善行，无论是大是小，在那一天都必被记在我们的账上。因此，即使你无法恢复到最严格的生活纪律，但至少若你能从目前的无度和放纵中稍作退避，这也并非不可能：只消立即着手，打开进入赛场的入口；但只要你还在外面徘徊，这事自然显得困难且不可行。因为在尝试之前，即使事情轻松易行，也往往会给我们带来很大的困难印象；但当我们实际投入尝试并冒险时，大部分痛苦就被消除了，信心取代了颤抖和绝望，减少了恐惧，增加了操作的便利，并使我们的美好希望更加强固。为此，恶者也把犹大拖出了这个世界，免得他有一个良好的开端，从而藉悔改回到他堕落前的状态。因为，虽然这话可能显得奇怪，但我仍不承认那罪大到超过了悔改带给我们的救助。因此，我祈祷并恳求你，从你灵魂中驱除这种撒殚式的思维方式，回到这个救恩的状态。因为如果我真的命令你一下子升到你从前的境界，你自然会抱怨这样做困难重重；但如果我现在只要求你起来，从那方向返回，你为什么犹豫、退缩、倒退呢？你难道没有见过那些在奢侈、醉酒、嬉戏及此生其他愚妄中死去的人吗？那些曾在大街上以众多排场和随从昂首阔步的人，如今在哪里？那些穿绸缎、满身香气、为奉承者设宴、常去戏院的人，如今在哪里？他们的一切排场如今变成了什么？一切都消失了——宴席的豪华花费、乐师的簇拥、谄媚者的殷勤、高声大笑、精神的松弛、心智的软弱、放荡、纵欲、挥霍的生活方式——一切都结束了。如今这一切都逃到哪里去了？那曾享受如此多关注和洁净的身体变成了什么？你到棺材那里去吧，看看尘土、灰烬、蛆虫，看看那地方的污秽，并痛苦哀号。但愿惩罚只限于灰烬！但现在，将你的思绪从棺材和这些蛆虫转到那不死的蛆虫、不灭的火、咬牙切齿、外面的黑暗、苦难与窘迫，以及拉匝禄和富翁的比喻上：那拥有如此多财富、身穿紫红袍的人，却连一滴水也得不到，尽管他处于如此巨大的需要之中。这个世界的事物，其本性并不比梦境更好。因为正如那些在矿中劳作或遭受其他更严厉刑罚的人，因着许多劳累和生活的极度痛苦而入睡，在梦中看见自己过着奢侈和繁荣的生活，醒来后并不感谢他们的梦；同样，那富翁在此生像在梦中一样变得富有，离世后便受到那痛苦的惩罚。考虑这些事，并以那火与你目前燃烧的欲望之火相对照，从这火炉中解脱自己吧。因为那在此地彻底熄灭这火炉的人，在另一个世界将不会经历那火；但若有人在此地未能战胜这火炉，他在离世后将被那火更猛烈地抓住。你希望今生的享受延长多久？我想你恐怕不会有超过五十年的时间活到极老，而且连这也完全不确定：因为连是否能活到晚上都无法确定的人，怎能信赖这许多年？不仅这不确定，还有我们境况变化的未知：因为常常生命虽延长了很长时间，奢侈的条件却并未随之延长，而是来了又匆匆离去。然而，如果你愿意，权且假设你能活那么多年，且不会遭遇任何命运逆转，这与无尽的岁月以及那些痛苦和难忍的惩罚相比，又算得了什么呢？因为在此地，善与恶都有终结，而且非常迅速；但在那里，两者都与不朽的岁月共存，其性质与现今的事物不可言喻地不同。\par

10. 当你听到火时，不要以为那世界的火像这火：因为这世界的火会烧毁并消灭它所抓住的一切；但那火却不断焚烧那些一旦被它抓住的人，永不停止：因此它也被称为不灭的。因为那些犯了罪的人也必要穿上不朽，不是为了荣耀，而是为了给那惩罚不断提供材料；这有多么可怕，言语无法描绘，但从微小事物的经验中，可以对那些大事略知一二。因为如果你曾进入一个加热过度的浴室，那么，我请你想想地狱之火；或者如果你因某种高烧而灼热，也将你的思绪转向那火焰，那么你就能清楚辨别其中的差别。因为如果一个浴室和一场发烧就这样折磨和困扰我们，那么当我们落入那环绕可畏审判座前的火河时，我们的处境将会如何？那时我们将因劳苦和难以忍受的痛苦而咬牙切齿；但没有人会来救助我们：是的，我们将大声呻吟，火焰更猛烈地加于我们，但我们只看到那些与我们一同受罚的人，以及极大的荒凉。谁能描述黑暗给我们的灵魂带来的恐惧？因为正如那火没有毁灭的能力，它也没有发光的能力：否则就不会有黑暗。那时这在我们心中产生的惊慌、战栗和极大的惊骇，只能在那一天充分体会。因为在那个世界，许多不同种类的折磨和惩罚的洪流从四面八方向灵魂涌来。如果有人问：\Quote{灵魂怎能承受如此众多的惩罚，并持续受罚无尽的岁月？}就让他想想这世上发生的事：有多少人曾忍受漫长而严重的疾病。如果他们死了，那不是因为灵魂被耗尽，而是因为身体衰竭；所以如果身体没有崩溃，灵魂就不会停止受苦。当我们接受了不朽不坏的身体时，就没有什么能阻止惩罚无限期地延续。因为在这里，确实不可能同时存在两种情况：即严厉的惩罚和持久的存在，二者互相冲突，因为身体的本性是脆弱的，不能同时承受两者；但一旦不朽的状态来临，这冲突就会结束，而这两种可怕的事物将有力地无限期地控制我们。所以，不要现在就这样设想，好像折磨的过度力量会毁灭灵魂：因为就连那时身体也不会经历这样的事，而是与灵魂一同存留，处于永恒的惩罚中，除了这之外，不再有任何终结。那么，多少奢侈、多少时光，你能用来与这惩罚和报复相权衡？你提出一百年或两倍的时间吗？与无尽的岁月相比，这算什么呢？因为如同一天之梦在整个一生中，此世之物与来世之境的对比也是如此。那么，有谁会为了看到一场好梦而选择永远受罚呢？谁如此愚昧，竟要寻求这种报应？因为我尚未指责奢侈，也尚未揭示其隐藏的苦涩：现在不是进行这些评论的时候，而是等你能摆脱它之时。因为现在，你被这情欲捆绑，若我说快乐是苦涩的，你会怀疑我胡言乱语；但当日，借着天主的恩宠，你从这疾病中释放后，你便会知道它的主题；另一时节，我现在要说的是：就算奢侈是奢侈，快乐是快乐，它们本身没有痛苦或可耻，我们将如何面对那为我们存留的惩罚呢？那时我们该怎么办，如果现在我们如同在影子和形象中享受快乐，而在那里却真实地经历永久的折磨，而我们本可在短时间内逃脱这些已提到的折磨，并享受为我们存留的美好事物？因为这也是天主慈爱的工作：我们的斗争不会拖延得很长，而是在短暂的、微小的眨眼之间（与另一世界相比，现世就是如此）之后，我们便为无尽的岁月领受胜利的冠冕。对那时受罚的灵魂来说，想到在这短暂的一生中本可补救一切，却忽略了机会，将自己交托给永恒的邪恶，这将是极大的痛苦。为了不让我们遭受这些，让我们在悦纳的时候，在救恩的日子，\BibleRef{格后 6:2}，在悔改力量强大的时候，振作起来。因为不仅已提到的邪恶，还有比这些更糟的其他邪恶在等待我们，如果我们懈怠。这些确实，以及比这些更痛苦的事，存在于地狱中；但失去美好事物所带来的痛苦、折磨和困窘是如此之大，即使没有为在此犯罪的人设定其他惩罚，仅此一项就足以比地狱中的折磨更痛苦地折磨我们，并扰乱我们的灵魂。\par

11. 请你思考来生的景况，尽可能思考；因为言语不足以充分描述；但根据所告诉我们的事，如同藉着某些谜语，让我们试着隐约看见它。我们读到：\BibleQuote{痛苦、悲伤和叹息已经逃走了。}\BibleRef{依 35:10}。那么，还有什么比这生命更蒙福呢？在那里不可能害怕贫穷和疾病；不可能看见任何人伤害或被伤害，激怒或被激怒，或愤怒，或嫉妒，或燃烧着任何无耻的情欲，或为生活必需品的供应而焦虑，或为自己失去某种尊严和权力而哀叹；因为我们内一切情欲的风暴都被平息并化为乌有，一切将处于和平、喜悦和欢乐的状态，一切宁静安详，一切是白昼和光明，不是这眼前的光，而是比这光更灿烂，正如这光胜过灯的光亮。因为在那世界，事物不被黑夜或云彩聚集所隐藏；身体在那里不被点燃焚烧；因为那里没有黑夜，没有黄昏，没有寒冷，没有炎热，也没有季节的任何变化；而是另一种方式，只有那些被认为配得上它的人才会知道；那里没有衰老，也没有衰老的任何邪恶，而与衰败有关的一切都被完全移除，不朽的荣耀充满各处。但比这一切更大的，是与基督、天使、总领天使和更高权能的永恒交往。现在请你观看天空，在思想中穿越到天空之外的区域，并思考整个受造物将要发生的变形；因为它不会继续像现在这样，而会远为辉煌美丽，如同金子比铅更闪耀一样，未来宇宙的状况将比现在更好；正如蒙福的保禄说：\BibleQuote{因为受造物本身也要从腐败的束缚中获得释放。}\BibleRef{罗 8:21}因为现在，既然它分享腐败，它经历许多这类身体自然经历的事；但那时，它摆脱了这一切，我们将看见它以不朽的形象展示它的美丽；因为它将要接受不朽的身体，它本身也将被改变为更高贵的状态。在那世界里，没有叛乱和争斗；因为圣人的群体有极大的和谐，彼此永远一致。在那里不可能害怕魔鬼、恶魔的阴谋、地狱的威胁或死亡，无论是现在的死亡，还是比这更糟的另一种死亡，而每一种这类恐惧都将被消除。就像某个王室的孩子，在卑微的装扮下被养大，并受到恐惧和威胁，以免因放纵而堕落，变得不配继承父亲的遗产，一旦他获得王室尊严，他立即将所有以前的衣服换成紫袍、冠冕和卫队，并充满信心地就位，将谦卑和服从的思想从灵魂中驱逐出去，并代之以其他思想；同样，那时对所有的圣人也将会是这样。\par

为了证明这些话不是空夸，让我们在思想中前往基督显圣容的山上；让我们看见祂如同在那里发光；然而即使那时，祂也没有向我们显示未来世界的全部光辉。因为那景象是为适应人眼，而不是现实的准确显现，从福音作者的话中就可明了。因为他说了什么？\BibleQuote{祂发光有如太阳。}\BibleRef{玛 17:2}但不朽身体的荣耀并不发出与这腐败身体相同的光，也不是那种可以忍受凡眼的光，而是需要不朽和不死的眼睛来观看。但那时在山上，祂向他们显示了在不对观看者的视力造成伤害的情况下他们所能看见的那么多；即便如此，他们也不能忍受，俯伏在地。请告诉我，如果有人把你领到一个明亮的地方，那里所有人都穿着金衣坐着，在人群中向你指出唯一一个穿着宝石衣、头戴冠冕的人，然后承诺把你安置在这群人中，你会不竭力去实现这个承诺吗？那么现在就在想象中睁开你的眼睛，看看那个集会，它不是由像我们这样的人组成，而是由比黄金和宝石、太阳的光线和所有可见光辉更有价值的人组成，而且不仅由人组成，还有比人更有尊严的存在——天使、总领天使、宝座、宰制者、率领者、掌权者。因为关于君王，甚至无法说出祂像什么；祂的美丽、恩宠、光辉、荣耀、伟大和庄严完全超越言语和思想。那么，我问，我们会为了逃避短暂的痛苦而剥夺自己这么大的福气吗？因为如果我们每天要忍受无数的死亡，甚至地狱本身，为的是看见基督在荣耀中来临，并被列入圣人的行列，难道我们不该承受这一切吗？请听蒙福的伯多禄说什么：\BibleQuote{我们在这里真好！}\BibleRef{玛 17:4}但如果他，当他看见未来事物的模糊影像时，立刻因那景象产生的喜悦而将一切其他事物从灵魂中抛开；那么当事物的实际现实呈现时，当宫殿打开，允许我们面对面地看到君王本人，不再模糊，或藉着镜子\BibleRef{格前 13:12}，而是面对面；不再藉着信德，而是藉着看见，那时人会说什么呢？\par

12. 确实，那些不甚明智的人大多以逃脱地狱为满足；但我要说，远比地狱更严厉的惩罚是被排斥于另一世界的荣耀之外，而且我认为，未能达到那荣耀的人，与其因地狱的苦楚而悲伤，不如因被拒于天堂之外而哀恸，因为仅此一项就比惩罚中其他一切更可怕。可是如今，当我们看到一位国王在众多护卫簇拥下进入王宫时，我们便认为那些亲近他、能分享他的言语和心意、并参与其余一切荣耀的人是有福的；即使我们拥有无数福分，我们也感受不到其中任何一样，而当我们目睹他周围之人的荣耀时，便自认不幸，尽管我们知道这种辉煌是易逝而不稳固的，或因战争、阴谋、嫉妒，或因其本身并不值得任何重视。然而，关系到万王之王时——祂所持有的不是大地的一隅，而是整个圆周，或者说，祂用手心容纳一切，用虎口测量诸天，用祂权能的话支撑万有，万民在祂面前如同无有，又如一滴唾液——对于这样一位君王，难道我们不会认为，不能列于祂周围的队伍中是最极端的惩罚，而仅仅满足于逃脱地狱？还有什么比这种灵魂状态更可怜呢？因为这位君王降临审判大地时，并非由一对白骡牵引，也不是乘坐金车，更不是身披紫袍、头戴冠冕。祂如何降临呢？请听先知大声呼喊，将此尽可能告知世人：一位说：\BibleQuote{天主必公开来临，我们的天主绝不缄默；烈火在祂面前燃起，狂风巨浪环绕祂。祂呼唤上天下地，为审判祂的子民。} 但依撒意亚描绘了我们面临的真正惩罚，如此说：\BibleQuote{看，上主的日子来临，残忍、愤怒、烈怒，使大地荒凉，消灭地上的罪人。因为天上的星辰、猎户座和整个天象都不发光，太阳在落下时昏暗，月亮不放射光芒；我必降灾祸于普世，惩罚恶人的罪孽，消灭目无法纪者的傲慢，压制骄傲者的骄矜；余下的人比未炼的金子更珍贵，人比蓝宝石更宝贵。因为上天震怒，大地因万军上主烈怒的暴怒在其降临时从根基动摇。}\BibleRef{依 13:9} 又说：\BibleQuote{窗户}他说\BibleQuote{要从天上打开，大地的根基要动摇；大地必全然混乱，大地必压低，极其困惑；大地必像醉汉和狂欢者一样猛烈摇晃；大地必如茅屋般颤抖，倒下不能再起，因为罪恶在其中盛行。在那一天，天主必在高天伸手攻击天上的万军，并攻击地上的列国；祂必将他们聚集在牢狱中，囚禁在堡垒里。} 玛拉基亚与此一致地说：\BibleQuote{看，全能的上主来临，祂来临的日子谁能忍受，祂显现时谁能站立？因为祂像炼金者的火，又像漂布者的碱；祂必坐下炼净、提纯，如同银子和金子。}\BibleRef{拉 3:2} 又说：\BibleQuote{看，}他说\BibleQuote{上主的日子来临，如同火炉，它要吞灭他们，所有外邦人和所有作恶的人必如碎秸，那来临的日子必点燃他们，万军的上主说；必不留下根或枝。}\BibleRef{拉 4:1} 那位大蒙宠爱的人说：\BibleQuote{我观看到宝座设立，万古常存者就坐，祂的衣服洁白如雪，头发纯净如羊毛；祂的宝座是火焰，轮子是燃烧的火；一条火河在祂面前奔流。千万事奉祂，万万侍立在祂面前。审判已设立，案卷已展开。} 过了一会儿，他说：\BibleQuote{我在夜间的异象中观看，见有一位像人子者，乘着天上的云而来，直到万古常存者那里，被引领到祂面前，祂被授予统治权、尊荣和国度，万民、各族、万舌都事奉祂。祂的统治是永远的统治，永不消逝，祂的国度必不毁灭。至于我达尼尔，我的灵在我里面战栗，我头脑中的异象使我烦乱。} 那时，天上的所有门户都打开了，或者更确切地说，天本身从中被除去，\BibleQuote{因为天，}我们读道，\BibleQuote{要像书卷一样被卷起，}\BibleRef{依 34:4} 像某帐篷的皮和罩子一样卷在中间，以便转化为更美好的形状。那时，万物充满了惊愕、恐惧和战栗；那时，连天使自己也被极大的恐惧所束缚，不仅是天使，还有总领天使、宝座、宰制、率领和掌权者。\BibleQuote{因为天上的众军，}我们读道\BibleQuote{都要动摇，}因为他们的同仆必须为今世的生活交账。\BibleRef{玛 24:29} 因为如果在今世，一个城市在统治者面前受审时，所有人都战栗，即使那些没有危险的人，那么当整个世界在这样一位法官面前受审时——祂不需要证人或证据，而是独立于这一切，将行为、言语和思想都呈现出来，如同图画一般，既展示给犯罪的人，也展示给不知情的人——难道一切权势不都自然惊慌失措吗？因为即使没有火河奔流，没有可怕的天使侍立在宝座旁，而只是人们被召集，一些人得赞美和钦佩，另一些人被羞辱地遣散，不得看见天主的荣耀（\BibleQuote{因为不虔敬的人，}我们读道，\BibleQuote{要被带走，免得看见上主的荣耀}），如果这是唯一的惩罚，那么失去这样的福乐岂不比一切地狱更刺痛被剥夺之人的灵魂吗？因为这种祸患有多大，如今无法用言语表达；但那时我们将在实际经验中清楚知道。但现在，我祈求你也在场景中加入惩罚，并想象人们不仅蒙羞、低头弯腰，而且被沿路拖向火中，被拉向酷刑器具，被交付给残忍的权势，并在所有行善者、行永生所应许之德者正被加冕、被宣布为胜利者、被引至王座前时，遭受这些痛苦。\par

13. 这些是那一日将要发生的事；但此后随之而来的事——与基督为伴的快乐、利益和喜乐——有什么语言能向我们描述呢？因为当灵魂回归到尊贵的正常状态，并能从此怀着极大的胆量瞻仰其主时，无法言说它从中获得何等大的快乐，何等大的利益，不只因实际拥有的美善而喜乐，更因确信这些美善永无止境。那一切喜乐无法用语言描述，也无法被理智掌握；但我将努力以模糊的方式，藉小事指示大事，使之显明。让我们细查那些在此生享受世界美善的人，我说的是财富、权力和荣耀，他们如何因快乐而欢腾，自认为不再在地上，尽管他们所享受的事物并不被承认为真正的美善，也不与他们长存，反而比梦境更快飞逝；即使它们持续片刻，其恩惠也只显现于今生范围之内，不能伴我们前行。如果这些事将拥有者提升到如此喜乐的程度，你认为那些被邀请享受天上无量福祉的灵魂，这些福祉永远稳固不动，其情境又如何呢？不仅如此，它们在数量和质量上都超越现世事物，达到人心中从未想过的高度。因为现今我们像子宫中的胎儿，被局限在狭窄的空间中，无法看见来世的光辉和自由；但当分娩的时刻到来，今生在审判之日将所有它所容纳的人诞生出来时，那些流产的从黑暗进入黑暗，从苦难进入更深的苦难；而那些完全成形并保存了君王肖像印记的，将被呈献给君王，并承担起天使和总领天使向万有之主所履行的职事。朋友啊，我恳求你，不要最终抹去这些印记，要速速恢复它们，并更完美地将它们铭刻在你的灵魂上。因为天主已将身体之美限制在本性的界限内，但灵魂的恩宠却摆脱了那原因带来的约束和束缚，因为它远高于任何身体的匀称；它完全依靠我们自己和天主的恩宠。我们的主，因着仁慈，特别以此方式尊崇我们的种族：祂将那些无大益、结局无关紧要的下等事物托付于本性的必然，而对于真正崇高的事物，祂使我们自己成为工匠。因为如果祂也将身体之美置于我们的控制之下，我们就会陷入过度的焦虑，将全部时间浪费在无益之事上，并严重忽视我们的灵魂。\par

因为，即使现在，当我们自身没有这种能力时，我们仍竭力努力，投身于涂脂抹粉，因不能真正产生身体的美，便用涂料、染料、发型、衣着、描眉及其他许多巧计狡猾地制造模仿：如果我们有能力将身体转变为真正匀称的形状，我们又能为灵魂和严肃事务留出多少闲暇？因为很可能，如果这是我们的营生，我们便不会有其他事，而将全部时间花在这上面：用无数的装饰打扮婢女，却让婢女的主母永远处于畸形和被忽视的状态。因此，天主将我们从这无益的劳碌中解救出来，把作用于更高贵元素的能力植入我们心中，那不能将丑陋的身体变为俊美的人，却能提升灵魂，即使它已处于丑陋的极点，也能达到恩宠的顶峰，使它如此可爱可欲，以致不仅善人渴想它，连万有的主宰、天主自己也不例外，就如圣咏作者论及这美丽时也说：\Quote{君王必恋慕你的美丽。}你们不也看见吗？在妓院里，那些丑陋无耻的妇女连拳击手、逃亡奴隶和角斗士也几乎不愿接受；但若有一个娟秀、出身高贵且端庄的妇女，因某种不幸被迫沦落至此，那么即使非常显赫伟大的人，也不会羞于与她成婚。如果人们竟有如此多的怜悯和对荣誉的蔑视，以致把那些常在妓院中受辱的妇女从束缚中释放，并使她们成为妻子，那么天主对那些因魔鬼篡夺而从原本高贵的境地坠入现世生活淫行的灵魂，更是如此。你会发现先知们充满了这样的例子，当他们向耶路撒冷讲话时；因为她行了淫，而且是一种新奇的形式，正如厄则克耳说：\Quote{所有的妓女都得到报酬，你却把报酬给了你的情人，你在所有女人中显出反常。}另一位又说：\BibleQuote{你像一只被遗弃的鸟，坐着等候他们。}\BibleRef{耶 3:2}这样行了淫的那一位，天主再次召叫她回来。因为所发生的充军，与其说是出于报复，不如说是为着悔改和改正；如果天主想要直接惩罚他们，就不会再带他们回家乡，也不会重建他们的城和圣殿，使之比先前更加辉煌：\BibleQuote{这殿最后的荣耀}他说\BibleQuote{要大过先前的。}\BibleRef{盖 2:10}如果天主没有将多次行淫的那一位排除在悔改之外，那么他更会拥抱你的灵魂，它如今是初次跌倒。因为确实没有贪爱肉体之美的人，即使他非常疯狂，会如此被对他情妇的爱火燃烧，像天主切望我们灵魂的得救一样；这一点我们可以从每天发生的事以及从圣经中看出。至少请看，无论在\BibleBook{耶肋米亚}{书}的开端，还是在先知们的许多其他地方，当他被轻视和藐视时，他是如何急忙前行，追求那些背离他的人的友谊；这他自己也在福音中说明了，说：\BibleQuote{耶路撒冷！耶路撒冷！你杀死先知，又用石头砸死那些派到你这里来的人；我多少次愿意聚集你的子女，好像母鸡把小鸡聚集在她的翅膀下，但你们却不愿意！}\BibleRef{玛 23:37}保禄写信给格林多人说：\BibleQuote{天主在基督内使世界与自己和好，不将他们的过犯归咎于他们，并将和好的话托付给我们。所以我们代基督作大使，好像天主借我们劝你们；我们替基督请求你们：与天主和好吧！}\BibleRef{格后 5:19}要想到这话如今是对我们说的。因为不仅缺乏信德，而且不洁的生活也足以产生这可憎的敌对。\BibleQuote{属肉体的心思}我们读到\BibleQuote{是与天主为敌。}\BibleRef{罗 8:7}那么，让我们打破这屏障，将它劈成碎片，完全摧毁，使我们能享受这有福的和好，使我们再次成为天主所宠爱的那一位。\par

14. 我知道你现在欣赏赫尔弥俄涅的恩宠，并认为世上没有什么能与她的容貌相比；但是，朋友，如果你愿意，你自己将在容貌和风姿上超过她，如同金像超过泥像。因为如果美貌出现在身体上时，尚且如此迷人和激动大多数人的心，那么当灵魂闪耀着它时，还有什么能与这种美丽和恩宠相比呢？因为这种身体之美的基础不过是痰、血、体液、胆汁和咀嚼过的食物流质。因为眼睛、脸颊和所有其他部位都是由这些东西滋润的；如果它们得不到这种滋润，皮肤日渐干枯，眼睛凹陷，整个面容的恩宠立即消失；因此，如果你考虑到那些美丽眼睛、笔直的鼻子、嘴巴和脸颊内部所储存的东西，你就会说这匀称的身体不过是粉饰的坟墓；里面的部分充满了如此多的污秽。此外，当你看到一块布上沾有这些东西，如痰或唾沫时，你连指尖都不愿触碰，甚至不忍心看它；然而你却对这些东西的仓库和贮藏室激动不已吗？但你的美不是这种，而是超越它，如同天高于地；不如说，它比这更好、更灿烂。因为没有人曾看到过灵魂本身，脱离身体；但即便如此，我将努力从另一源头向你展示这灵魂的美。我是说从更大能力的例子。请听这些能力的美如何打动了那位深受爱戴的人；因为想要描述它们的美，却找不到同样性质的身体，他只得求助于金属物质，甚至对此也不满足，而是借用了闪电的光芒来作比喻。\BibleRef{达 10:6} 如今，如果那些能力即使没有透露其纯粹赤裸的本质，而只以非常模糊和影子的方式显现，却仍如此明亮，那么当它们摆脱一切遮蔽时，其本来面目该是什么样的呢？现在我们应该形成某种关于灵魂之美的形象。我们读到：\BibleQuote{因为他们将成为与天使相等的。} \BibleRef{路 20:36} 至于身体，较轻、较精细、并趋向非物质的那些种类，远比其他更好更奇妙。至少天空比大地更美，火比水更美，星辰比宝石更美；我们欣赏彩虹远胜于紫罗兰、玫瑰和地上所有的花朵。简言之，如果有可能用肉眼看到灵魂之美，你就会嘲笑这些身体的比喻，因为它们如此微弱地向我们呈现了灵魂的优美。那么，我们不要忽视这样的财宝，也不要忽视如此大的幸福，尤其是当我们对那种美的接近因我们对未来之事的希望而变得容易时。我们读到：\BibleQuote{因为我们这短暂轻微的苦难，正为我们成就极重无比永远的荣耀；我们不是顾念所见的，乃是顾念所不见的；因为所见的是暂时的，所不见的是永远的。} \BibleRef{格后 4:17} 如果蒙福的保禄称你所知道的这些苦难为轻微容易的，因为他没有顾念所见的事物，那么仅仅停止放荡就容易得多了。因为我们没有召唤你去经歷他所遭遇的那些危险，也没有召唤你去受他每日所犯的死亡，那经常的鞭打和鞭笞，锁链，全世界的敌视，本族人的仇恨，频繁的守夜，长途旅行，船难，强盗的攻击，亲族的阴谋，因朋友而受的困苦，饥饿，寒冷，赤身露体，火烧，因属于他和不属于他的人而沮丧。我们现在不要求你做任何这些事；我们所要求的只是你把自己从那可咒诅的奴役中释放出来，回到你以前的自由，考虑你放荡带来的惩罚和你以前生活方式带来的尊荣。因为不信者对复活的思考如此冷淡，从不害怕这类事情，这并不奇怪；但我们这些对来世之事比现世之事更确信的人，却以这种悲惨和可怜的方式生活，对这些事情的记忆毫无感动，反而陷入极端麻木的状态——这是极其不合理的。因为当我们这些相信的人做出不信者的事，或者甚至比他们处境更糟时（因为在他们中间有些人因生活美德而卓越），还有什么安慰、什么借口留给我们呢？许多商人遭遇船难后并没有放弃，而是继续同一条航程，而且他们遭遇的损失并非由于自己的疏忽，而是由于狂风的力量；而我们这些对结局有理由充满信心，并且确知如果我们不愿意，无论是船难还是任何意外都不会给我们带来损害的人，难道不重新开始工作，像从前一样经营我们的生意，反而闲懒不动，束起双手吗？但愿我们只是把双手\Emph{束起}，而不是用它们\Emph{对付}自己，那是完全疯狂的标志。因为如果任何一个拳击手，离开他的对手，反手打自己的头，击打自己的脸，我们难道不把他列在疯子中吗？因为魔鬼已经把我们打翻在地，所以我们应该站起来，而不是再被拖下去，自行跌倒，在魔鬼的打击之上再加自己的打击。因为蒙福的达味也曾经历过像你现在所遭遇的这样的跌倒；不仅如此，还有另一个跌倒随之而来。我指的是谋杀的跌倒。那又怎样？他保持躺卧了吗？他不是立即精力充沛地站起来，准备好与敌人作战了吗？事实上，他与敌人搏斗得如此勇敢，以至于在他死后，他仍是其后裔的保护者。因为当撒罗满犯了极大的罪，该当无数次的死亡时，天主说祂会为他保留完整的王国，这样说：\BibleQuote{我必将王国从你手中夺去，赐给你的臣仆。然而，我不在你活着的时候这样做。} 为什么？\BibleQuote{为了你父亲达味的缘故，我要从你儿子手中夺去。} \BibleRef{列上 11:11} 再者，当希则克雅即将面临最大的危险时，尽管他是个义人，天主却说为了这位圣人的缘故要帮助他。祂说：\BibleQuote{我要用我的盾牌掩护这城，为我自己和我仆人达味的缘故拯救它。} \BibleRef{列下 19:34} 这就是忏悔的力量。但如果他自己决定，像你现在这样，从此以后不可能取悦天主，如果他心里说：\Quote{天主以极大的尊荣尊崇了我，给了我一个先知的位置，将管理同胞的权柄托付给我，从无数危险中拯救了我，那么，在我受了这样大的恩惠之后得罪了祂，犯了最严重的罪行，我还能恢复祂的恩宠吗？} 如果他这样想，他不仅不会做后来所做的事，而且会加重他以前的恶行。\par

15. 因为不仅身体的创伤若被忽视会导致死亡，灵魂的创伤也是如此。然而我们竟愚昧到如此地步：对前者倍加呵护，对后者却置之不理。虽然在身体方面，许多创伤常常天然是不可治愈的，但我们并不放弃希望，甚至在听到医生不断宣告\Quote{此病无法用药物根除}时，我们仍恳求他们至少设法减轻一些痛苦。但在灵魂方面，却没有不治之症，因为它不受自然必然性的支配。在这里，我们却仿佛病症陌生一般，疏忽大意，陷于绝望。而在疾病的本性可能使我们陷入绝望的地方，我们却仿佛极有康复希望那样费尽心力；但在本不应放弃希望的地方，我们却停止努力，变得漫不经心，如同事已无望。由此可见，我们对身体的重视远超过灵魂。这也正是我们连身体也无法保全的原因。因为忽视主导要素，却将全部热忱倾注于次要事物的人，最终会丧失两者；而遵守正确秩序、保存并珍爱权威要素的人，即使忽略次要要素，也会因保全首要要素而保存它。基督也曾以此晓谕我们，祂说：\Quote{不要害怕那些杀害身体，却不能杀害灵魂的；但更要害怕那能把灵魂和身体都毁灭在地狱里的。}\BibleRef{玛 10:28}\par

好吧，我是否说服了你，人绝不该对灵魂的紊乱绝望，视为不治？还是我必须再提出其他论据？因为即便你对自己绝望千万次，我也决不会对你绝望，我也决不会亲自犯下我指责他人的那罪。然而，一个人对自己绝望与别人对他绝望并不是一回事。因为在别人身上怀有这种猜疑的人，或许容易获得宽恕；但自己怀有这种猜疑的人，则不会。这究竟是为什么？因为前者对别人的热忱和补赎没有控制权，而人对自己的热忱和补赎则拥有唯一的主权。不过即便如此，我也不会对你绝望；即使你无数次这样受影响，因为或许会有回归德行、恢复你先前生活方式的一天。现在请听以下的：尼尼微人听见先知大声宣布，明确威胁说：\BibleQuote{再过三天，尼尼微就要倾覆。}那时他们并没有灰心，反而，尽管他们不确信自己能够感动天主的心意，或者更确切地说，从神谕中理应怀疑相反的结果（因为话语没有任何条件，只是简单的宣告），那时他们仍表现出补赎，说：\BibleQuote{谁知道天主是否会后悔、回心转意，撤回祂的烈怒，使我们不至灭亡？天主看见他们所行的事，即他们离开了邪恶的行径，天主就后悔了祂曾说要将降在他们身上的灾祸，并没有实施。}\BibleRef{纳 3:9}如今如果野蛮、无理性的人能领会这么多，那么被神圣教义培养、并在历史和实际经验中见过如此多此类榜样的人，我们更应该这样做。\BibleQuote{因为我的计谋}我们读到\BibleQuote{不是你们的计谋，我的道路也不是你们的道路；天离地有多高，我的想法和你们的想法、我的计谋和你们的计谋之间的距离就有多远。}如果我们常常因家奴冒犯我们，当他们承诺变得更好时，就接纳他们，使他们回到原先的份额，甚至有时赐予他们比从前更大的说话自由，那么天主岂不更是如此？因为如果天主创造我们是为了惩罚我们，你很可能绝望，并质疑自己获得救恩的可能性；但若祂创造我们无他，只出于祂的善意，并为了让我们享受永久的福乐，并且祂从第一天直到现在，所做所谋一切都为此目的，那么有什么能使你怀疑呢？我们是否曾严重激怒祂，甚至超过任何人？这正是我们尤其应当放弃当前行为、为过去补赎并展现巨大改变的原因。因为我们曾犯下的恶行，并不及我们将来不愿改变更能激怒祂。犯罪不过是人性的软弱，但继续犯同样的罪就不再属于人性，而完全是魔鬼的行径。请看天主如何通过先知的口更责备后者而非前者。\BibleQuote{因为}我们读到\BibleQuote{我对她说，在她行了这一切淫乱的事之后，归向我吧，她却没有归来。}\BibleRef{耶 3:7}又从另一处：当祂想表明自己对我们的救恩怀有多大的渴望时，祂听见百姓在多次过犯后承诺行走正路，就说：\BibleQuote{谁使他们有这样的心，好敬畏我，日日遵守我的诫命，使他们和他们的子孙永世得福？}\BibleRef{申 5:29}梅瑟与他们辩论时说：\BibleQuote{以色列啊，现今上主你的天主向你要求什么？不过是敬畏上主你的天主，遵行祂的一切道路，爱祂。}\BibleRef{申 10:12}那么，那位如此渴望被我们爱、并为达此目标而做一切、且因对我们的爱不惜牺牲自己独生子、并且认为我们一旦与祂和好是可喜之事的天主，祂在我们补赎时岂不会欢迎并爱我们？请听祂通过先知口所说的：\BibleQuote{先陈明你的过犯，好使你成义。}\BibleRef{依 43:26}祂这样要求我们，是为了加深我们对祂的爱。因为当一位爱者，在遭受所爱者多次侮辱后，仍不熄灭对他们的喜爱，他之所以费心公开那些侮辱，唯一的原因就是通过显示他情感的力量，使他们感受到更大更热烈的爱。既然罪的告解带来如此大的安慰，那么努力通过行为洗去罪恶就更当如此。因为若非如此，那些曾偏离正道的人不允许再归回，那么或许除了少数几个屈指可数的人外，没有人能进入天国；但事实上，我们将发现最杰出的人中许多是曾跌倒过的。因为那些在恶事上表现极其猛烈的人，也会转而同样表现在善事上，因为他们意识到自己欠下了多大的债；基督在论及那妇人时对西满所说的话也表明了这一点：\BibleQuote{你看见吗？}祂说，\BibleQuote{这妇人？我进了你的家，你没有给我水洗脚，她却用眼泪洗了我的脚，又用头发擦干。你没有给我亲吻，她却从我进来起不住地亲吻我的脚。你没有用油抹我的头，她却用香膏抹我的脚。所以我告诉你：她的许多罪得了赦免，因为她爱得多；但那赦免少的，爱得就少。就对她说：你的罪赦了。}\BibleRef{路 7:44}\par

16. 为此，魔鬼也知道，那些犯了重大罪行的人，一旦开始悔改，便因意识到自己的过犯而充满热忱；魔鬼害怕战栗，唯恐他们着手这项工作；因为一旦开始，他们便不再能被阻止，在悔改的影响下如同火焰燃起，使自己的灵魂比纯金更纯净，被良心和往昔罪恶的记忆如同强劲的风催向德行的港湾。正是在这一点上，他们比那些从未跌倒的人更有优势：他们焕发出更猛烈的活力；只要如我所说，他们能把握住开端。因为艰难而难以成就的任务，在于能够踏上入口，到达悔改的门廊，并在那里击退和推翻正在猛烈狂怒攻击我们的敌人。但此后，一旦他被打败，在他曾强大的地方跌倒，他便不会如此狂暴，我们将获得更大的力量，并更轻松地跑完这场善赛。因此，让我们今后着手回归，让我们赶紧奔向天上的城，在那里我们已被登记，在那里我们也已被指定为公民安居。因为对自己绝望不仅带来这恶果：它关闭了那座城的大门，驱使我们陷入更大的懈怠和轻慢，而且还将我们投入撒殚式的鲁莽。因为魔鬼之所以变成如今这个样子，唯一的原因就是他首先绝望，然后从绝望陷入鲁莽。因为灵魂一旦放弃自己的救恩，便不再察觉自己正在下沉，选择做和说一切有损于自己救恩的事。正如疯子一旦脱离健康状态，便对任何事物不再恐惧或羞愧，无所畏惧地敢于一切，即使必须落入火中、深水或悬崖；同样，那些被绝望的狂乱所攫取的人从此便无法驾驭，冲入各种恶行，若死亡不来终止这种疯狂和激烈，他们便对自己造成无限伤害。因此，在我恳求你，在你深深陷入这种醉酒之前，恢复理智，振作起来，摆脱这种撒殚式的发作，温柔而逐渐地进行，若不能一次完成的话。因为对我来说，似乎更容易的方式是一劳永逸地挣脱所有束缚你的绳索，转而进入悔改的学校。但如果这对你来说似乎困难，即你愿意踏上通向更好事物的道路，那就单单踏上它，抓住永生。是的，我以你从前的声望、以你曾经拥有的那份自信，恳求并哀求我们再次看见你站在德行的顶峰，保持与从前一样的坚持。请宽恕那些因你而绊倒、跌倒、变得更为懈怠、对德行之道绝望的人。因为沮丧如今占据了弟兄们的队伍，而不信者及那些倾向于懈怠的青年人中却充满快乐和喜悦。但如果你重归从前严谨的生活，结果将反过来：所有我们的羞耻将转移到他们身上，而我们将享有极大的自信，看到你再次加冕，并以比从前更辉煌的方式宣告胜利。因为这样的胜利带来更大的荣誉和喜乐。你不仅将获得自己成就的赏报，还将获得对他人的劝勉和安慰，被树立为一个显著的榜样——若有人陷入同样境况，便可受鼓励站起并恢复。不要忽视如此有益的机会，也不要把我们的灵魂拖入阴间的悲伤中，而是让我们重新呼吸，抖落因你而压在我们身上的沮丧阴云。因为如今，我们忽略了自己的烦恼，哀悼你的灾祸；但如果你愿意清醒过来，看清一切，加入天使的行列，你将使我们脱离这悲伤，并除去大部分的罪恶。因为那些悔改后回转的人能闪耀出极大的光辉，往往甚至超过从未跌倒的人，我已从神圣著作中证明了这一点。至少，税吏和娼妓就这样继承天国，许多后来的居于先前的之前。\par

17. 但我还要告诉你一些发生在我们这个时代、你自己也可能亲眼目睹的事件。你大概认识那位年轻的腓尼基人，乌尔巴诺之子，他不幸过早成为孤儿，却拥有大量钱财、许多奴隶和土地。此人起初彻底告别了学校的学业，脱下他从前穿着的华丽衣裳和一切世俗的荣华，突然披上一件破旧的外衣，退隐到山野的孤寂之中，展现了高度的基督哲学，不仅与他的年龄相称，而且是任何伟大非凡之人都可能展现的。此后，他被认为配得领受神圣奥迹的入门，在德行上取得了更大的进步。众人都欢喜雀跃，光荣天主，因为一个在财富中长大、有显赫祖先、尚且年少的人，竟能突然践踏此生的一切虚荣，升到真正的高处。当他处于这种状态并受人赞叹时，某些腐败的人——按亲属关系监管着他——又把他拖回从前世俗的海洋。于是，他抛开所有习惯，再次从山上下到广场中央，骑着马，带着一大群随从，遍行全城；他不再愿意过节制的生活，因为被过度的奢华所点燃，被迫陷入愚蠢的爱情纠葛，所有与他交往的人没有不绝望于他的救恩的；他被如此众多的谄媚者包围，加上孤儿、青春和巨大财富的陷阱。那些凡事吹毛求疵的人指责最初引导他走上这条生活道路的人，说他既失去了属灵目标，又不再能在管理自己的事上有任何用处，因为他过早放弃了学业，从而无法从中获得任何益处。当这些话被说开，人们深感羞愧时，某些圣善的人——他们曾多次在这种追猎中成功，并凭经验完全知晓那些怀有对天主希望的人不应对此类人物绝望——持续地监视着他，每当他出现在市场时，他们就上前问候他。起初，他骑在马上，侧眼瞥视着他们跟随在旁；如此的厚颜无耻最初占据了他。但他们仁慈而仁爱的人，完全不因这种对待而感到羞耻，只持续注意一件事：如何将羔羊从狼群中救出；事实上他们靠着坚持到底确实做到了。因为后来，仿佛被某种突然的打击所转变，被他们极大的关切所羞愧，每当他在远处看见他们走近，就立刻下马，弯腰俯身，默默聆听他们口中所出的一切，久而久之，他甚至对他们表现出更大的尊敬和敬意。然后，藉着天主的恩宠，逐渐将他从这一切缠累中解救出来，他们又将他交还给从前那种隐居和虔诚的默观生活。如今他变得如此杰出，以致他从前的生活似乎与他堕落之后的生活相比不足挂齿。因为他亲身体验了陷阱，将全部财富施舍给穷人，解除了自己对这类事务的一切牵挂，从而斩断了那些欲设谋害他之人的一切攻击借口；如今他踏上了通往天堂的道路，已经抵达了德行的终点。\par

这个人确实在年轻时跌倒又站起；但另一个人，在沙漠中寄居时忍受了巨大的劳苦，仅有一位同伴，过着天使般的生活，且已渐入老境，却因某种撒殚般的心态和疏忽，给了恶者一个我不知道的小小漏洞；虽然他自从度隐修生活以来从未见过女人，却陷入了与女人交往的强烈欲望中。首先，他恳求同伴供给他肉和酒，并威胁说如果得不到，他就要下到市场去。他说这话，与其说是出于对肉的渴望，不如说是想找一个进入城市的借口和托词。同伴对此感到困惑，害怕如果阻止他，可能会把他逼入某种大恶，便任由他满足了这个渴望。但当同伴意识到这是一个陈腐的计谋时，他就公开抛弃了羞耻，揭穿了他的伪装，说自己必须亲自下到城市去；同伴无力阻止他，最终放弃了努力，远远地跟着他，观察这次返回可能意味着什么。看见他进入了一家妓院，并知道他在那里与一名娼妓行淫，同伴便等他满足了那污秽的欲望后，在他出来时，举起双手迎接他，拥抱并热切地亲吻他，对发生的事没有说出任何责备，只是恳求他，既然已满足了欲望，就再回到荒野中的住所去。这人被同伴极大的宽容所羞愧，立刻因自己所行的事而心痛悔恨，跟随他上了山；在那里他恳求那人把他关进另一间小屋，关闭住所的门，在特定的日子供应他面包和水，并告诉询问他的人说他已安息。他说了这话并说服了同伴后，便将自己关在里面，不断以禁食、祈祷和眼泪，从灵魂上擦去他罪恶的污点。不久之后，当邻近地区遭遇干旱，当地所有人都为此哀叹时，有一个人在异象中被命令出发，劝诫这位隐修士祈祷，以结束干旱。他出发了，带着同伴，发现从前与他同住的人独自在那里；询问另一个人时，被告知他已死了。但他们相信自己受了骗，便再次祈祷，又在同一异象中听到先前听到的话。于是他们围住那个确实欺骗了他们的人，恳求他把另一个人指给他们看；因为他们宣称他没有死，而是活着的。听到这话，他意识到他们的约定暴露了，便带他们到那位圣人那里；他们打破了墙壁（因为他甚至堵住了入口），全都进去，俯伏在他脚前，告诉他发生的事，恳求他帮助他们抵抗饥荒。他起初抗拒，说自己远远没有那样的信心；因为他总是将罪摆在他眼前，仿佛刚刚发生一样；但当他们把所发生的一切都告诉他时，他们才说服他祈祷；他祈祷后结束了干旱。至于那个起初是\Person{载伯德}的儿子\Person{若望}的门徒，但后来长期成为强盗头目，然后又因蒙福宗徒的圣手被俘，从强盗的巢穴和洞穴回到先前美德的年轻人，他所遭遇的事，你们并非不知，而是像我一样详细知道：我也曾多次听你们钦佩那位圣人的极大谦逊，他如何首先亲吻那年轻人沾满血的手，拥抱他，从而使他回到先前的状态。\par

18. 此外，圣\Person{保禄}不仅接纳了那无用的、逃亡的窃贼\Person{敖乃息摩}，因他已悔改，而且还请求他的主人以与自己的老师同等的尊敬来对待这个悔改者，这样说：\Quote{我为了我在缧绁中所生的儿子\Person{敖乃息摩}求你，他从前于你无益，但如今于你和我的确有益。我现在打发他到你这里来，你收下他，他就是我的心肝；我本来愿意将他留在我这里，叫他在你的地位上替你在福音的锁链中伺候你；但未得你的同意，我决不愿意做什么，叫你的善行不是出于勉强，而是出于甘愿。因为他暂时离开你，也许正是为叫你永远收下他，不再如同一个奴隶，而是高于奴隶，是一个可爱的弟兄：尤其为我是如此，何况为你呢？不拘在肉身方面，或是在主内方面，更是如此。所以，若你以我为同伴，就收纳他如同收纳我一样。} 同一宗徒在致\Place{格林多}人书函中又说：\Quote{免得我来到时，我要为那些从前犯罪而没有悔改的许多人悲伤。} 又说：\Quote{我从前说过，现在再说一次：我若再去，决不会宽容。} 你看到了吗？他为之悲伤和不肯宽容的是谁？不是那些犯罪的人，而是那些没有悔改的人；而且不单纯是没有悔改的人，而是那些被一再召唤去做这事却不肯受劝的人。因为\Quote{我从前说过，现在又事先说明，好像第二次在场，如今不在的时候写信} 这句话，恰恰预示着我们害怕如今在我们身上可能发生的事。因为虽然当时威胁\Place{格林多}人的\Person{保禄}不在场，但当时藉着他的口说话的\Person{基督}却在场；如果我们顽固不化，祂就不会宽容我们，却要在此世和来世重重打击我们。\Quote{让我们藉着承认己罪来迎接祂的威仪罢，} 让我们在祂面前倾吐我们的心声。因为经上记载：\Quote{你犯了罪，不要再加添，并为你从前的罪过祈祷。}\BibleRef{德 21:1} 又说：\Quote{义人在起初就是自己的控告者。} 我们不要再等待控告者，而要先占据他的位置，藉着我们的坦诚使我们的判官更为仁慈。现在我知道你确实承认你的罪，并称自己可怜至极；但这不是我唯一希望的，我渴望你相信这能使你成义。因为只要你使这承认毫无益处，即使你控告自己，也将无法停止由此产生的诸罪。因为除非一个人首先说服自己这样做是有益的，否则没有人能带着热忱和正确的方法做成任何事。就连农夫，撒下种子后，除非他指望收获，绝不会收割。因为谁愿意白白劳累，如果劳作中得不到任何益处？同样，播种言语、眼泪和认罪的人，除非怀着美好的希望去做，否则仍会被绝望的邪恶所压制而无法停止犯罪；正如那对任何果实收成绝望的农夫，不会再阻止那些损害种子的东西一样，那流着眼泪播种认罪却不指望由此获得任何益处的人，也无法推翻那些破坏悔改的东西。而破坏悔改的事，就是再一次陷入同样的邪恶之中。经上记载：\Quote{因为有一人}我们读到，\Quote{谁建造、谁拆毁，除了劳苦之外，他们得了什么？那因接触死尸而浸在水中，然后又去触碰它的人，他的水洗了什么益处？} 同样，如果有人为了罪而禁食，然后又走开去做同样的事，谁又能垂听他的祈祷？经上又说：\Quote{人若从义德转向罪恶，上主必为他准备刀剑。}\BibleRef{德 26:28} 并说：\Quote{人如犬转回自己的呕吐物，成为可憎的，愚昧人因邪恶转回自己的罪过。}\BibleRef{箴 26:11}\par

19. 那么，不要仅仅自我控告，陈述你的罪过，而要像一个应当通过悔改的方法成义的人那样行事；因为这样你就能使你那告解的灵魂感到羞愧，使它不再陷入同样的罪过之中。因为严厉地控告自己、称自己为罪人，这可谓连不信者也常做。至少，许多属于舞台的人，包括男人和女人，他们惯常行最大的无耻，也称自己为可怜者，但并非出于正确的目的。因此，我甚至不称此为告解；因为他们宣扬自己的罪过，并非伴随着灵魂的痛悔，也不是苦泪，更不是生活的皈依；事实上，他们中有些人这样做是为了在听众面前赢得直言不讳的名声。因为当别人宣布过失时，似乎不如犯罪者自己报告时那样严重。而那些在强烈绝望影响下陷入麻木状态、藐视同辈意见的人，以极大的厚颜无耻宣告自己的恶行，仿佛那是别人的行为。但我不愿你成为这些人中的任何一个，也不愿你从绝望中被带出而走向告解，而是带着美好的期望，在砍除绝望的整个根源之后，向相反的方向表现出热忱。这绝望的根源和母亲是什么？是懒惰；或者更准确地说，人们不仅称它为根源，还称它为保姆和母亲。因为正如羊毛腐烂滋生蛀虫，而蛀虫又反过来增加腐烂；同样，在这里，懒惰滋生绝望，而绝望又反过来滋养懒惰；这样，它们以这该死的交换互相供给，获得了不小的额外力量。如果有人砍断其中之一，将其击碎，他就容易制服剩下的一个。因为一方面，不懒惰的人永远不会陷入绝望；另一方面，被美好希望支持、不绝望于自身的人，也不会陷入懒惰。那么，请撕裂这对搭档，将轭击碎，我指的是那种多变而令人沮丧的思维习惯；因为将这两者结合在一起的东西不是单一的，而是在羞耻和性质上多样的。这又是什么呢？那就是一个悔改的人做了许多伟大而良善的事，但与此同时，他犯下了与那些善事相当的罪过，而这尤其足以将他推入绝望，仿佛已经建立起来的建筑全部倒塌，他所付出的所有劳动都归于徒劳和虚无。但这一点必须考虑到，这样的推理必须被击退，因为如果我们不提前储存与随后犯下的罪过相当的善工，那就无法阻止我们严重而彻底地沉沦。然而实际上，（正当行为）就像一件坚固的胸甲，不让尖锐而苦涩的箭矢完成其工作，即使它自己被刺穿，也避免了身体遭受许多危险。因为带着许多善行和恶行离世的人，在那里的惩罚和痛苦中会得到某种缓解；但如果一个人缺乏这些善工，只带着恶行离去，那么当他被引向永罚时，他将遭受多么巨大的痛苦，这是无法言说的。因为在那里，恶行与非恶行之间将有一个天平；如果后者压过秤盘，它们将在很大程度上拯救行为者，恶行所带来的伤害不足以将他从首位拖下；但如果恶行超过，它们就将他带入地狱之火，因为他的善行数量不足以抵挡这猛烈的冲击。这些事不仅是由我们自己的推理所提示，也是由神圣的谕示所宣告的；因为祂亲自说：\BibleQuote{祂要按每人的行为予以赏报。}\BibleRef{罗 2:6} 不仅在地狱中，在天国里也会发现许多差别；因为祂说：\BibleQuote{在我父的家里有许多住处；}\BibleRef{若 14:2} 并且：\BibleQuote{太阳的光辉是一样，月亮的光辉是另一样。}\BibleRef{格前 15:41} 如果祂在处理如此重大的事情时如此精确地说话，看到祂宣称即使在另一个世界里，一颗星与另一颗星之间也有差别，这有什么奇怪呢？既然知道这一切，让我们永不停止行善，也不厌倦；如果我们不能达到太阳或月亮的等级，也不要轻视星星的等级。因为只要我们至少显示出如此程度的德行，我们就能在天堂中有位置。即使我们可能没有成为金子或宝石，只要我们占据了银子的等级，我们就能留在根基中；只是不要让我们再回到那容易被火吞噬的材料中去，也不要因为我们无法完成大事就放弃小事，因为那是极端愚昧的一部分，我相信我们不会经历。因为正如物质财富的增长，如果爱财者不轻视最小的收益，那么精神财富也是如此。因为奇怪的是，审判者不会忽略甚至一杯凉水的报酬，而我们如果成就并非全然伟大，却忽视做小事。因为不轻视小事的人，也会对大事十分热忱；而忽视前者的人，也会远离后者；为防止这种情况发生，基督甚至为这些小事规定了巨大的赏报。因为有什么比探访病人更容易呢？然而祂甚至用巨大的赏报来回报这事。那么，把握永生，在主内喜乐，并恳求祂；再次负起轻松的轭，屈就于轻省的重担，结束你与开始相称的生命；不要让自己错过如此巨大的财富之流。因为如果你继续以行为激怒天主，你将毁灭自己；但如果在此之前，在很少的损害发生之前，在你所有的耕作都被洪水淹没之前，你堵塞了邪恶的渠道，你将能够恢复被破坏的东西，并再增添不少出产。考虑到这一切，抖掉灰尘，从地上起来，你对敌人将变得可畏；因为他自己确实已推倒了你，仿佛你永远不会再起来；但如果他看到你再次举起手来攻击他，他将受到如此意外的打击，以至于他将不再那么急于再次把你推倒，而你自己将来也会更安全，免受此类伤害。因为如果别人的灾难足以教训我们，那么我们自己所遭受的就更足以教训我们了。这就是我期望很快在你亲爱的自身身上看到的事，并且藉着天主的恩宠，你再次变得比以前更加光辉，并显示出如此伟大的德行，以至于你在上界成为他人的保护者。只是不要绝望，不要退缩；因为在我每一句言语中，无论我在何处见到你，以及藉着别人的口，我将不会停止重复这一点；如果你听从这一点，你将不再需要其他药物。\par

% structure: p0027 source=h2 target=section reason=relative-heading-rank; base=h2

\section{第二封信}

一、倘若能用文字表达眼泪和呻吟，我定会将现时寄给你的这封信充满它们。如今我哭泣，并非因为你为你的家产焦虑，而是因为你已从弟兄的名册中涂抹了自己的名字，因为你已践踏了你与基督所立的盟约。这是我战栗的缘由，这是我痛苦的起因。为此我恐惧战兢，深知背弃这盟约将为那些投身于这场崇高争战、却因懈怠而擅离本位者带来严厉的定罪。而这样的人所受的惩罚比其他人更重，理由显而易见。因为无人会控告一个平民逃避兵役；但一个人一旦成为士兵，若被发现擅离队伍，他将冒着承受最严厉刑罚的危险。亲爱的\Person{德奥多罗斯}，摔跤手跌倒不足为奇，但若他继续倒在地上才是奇怪；同样，战士受伤并非可悲之事，但若在受伤后绝望，忽视伤口，才是可悲的。没有一位商人因一次海难和货物损失就停止航行，而是再次穿越大海、波涛和浩瀚的洋面，重获过去的财富。我们也看到竞技者，在多次跌倒后仍在胜利的冠冕中得胜。并且，从前时常有士兵曾逃跑，后来却成为勇士，战胜了敌人。许多因酷刑压迫而否认基督的人，后来再次争战，最终额戴殉道者的冠冕离去。但如果他们中的每一个在第一次打击后就绝望，就不能收获后来的益处。如今也是如此，亲爱的\Person{德奥多罗斯}，因为仇敌使你稍微动摇，你不要把自己再推入深渊，而要勇敢地站起来，迅速返回你所离开的地方，不要把这短暂的打击视为羞辱。因为如果你看见一个士兵从战场受伤归来，你不会羞辱他；因为抛弃武器、远离敌人才是羞辱；只要一个人仍在战斗，即使他受伤并短暂后退，没有人会如此冷酷或对战争缺乏经验，以致责备他。避免受伤是非战斗人员的事；而那些以极大勇气进攻敌人的人，有时可能会受伤和失败；这正是现在发生在你身上的事；因为你突然试图摧毁那蛇，却被咬了一口。但要鼓起勇气，你只需稍加警惕，这伤口的痕迹将不复存在；更确切地说，因天主的恩宠，你将粉碎那恶者自身的头；也不要因为你在起头就很快受阻而烦恼。因为那恶者锐利的眼睛已经看穿了你灵魂的卓越，并从许多迹象中猜出将有一个勇敢的对手对他日益强大；他预料到一个如此带着巨大热忱迅速攻击他的人，若持之以恒，必能轻易克服他。因此他勤勉、警醒，并大力鼓动起来反对你，更确切地说，是反对他自己的头，只要你勇敢地站稳。因为谁不惊叹你迅速、真诚而热切地向善的转变？因为对于美味珍馐你已漠不关心，锦衣华服已被轻视，一切虚浮的排场都被摒弃，而对世俗智慧的一切热忱突然转移到神圣的典籍上；全日用于阅读，全夜用于祈祷；不再提家族的尊贵，也不思念你的财富；而拥抱弟兄的膝、急忙跑到他们脚前，你认识到这比出身豪门更高贵。这些事激怒了那恶者，这些事激起他更激烈的争斗；但他并未给予致命的一击。因为若在长时间、持续禁食、睡硬地以及其余操练之后，他击败了你，即使如此也不需绝望；但人们会说，若在多番劳苦、辛劳和胜利之后失败，损失就大了；然而既然你一与他交手他就击倒你，他所成就的只是使你更加渴望与他争战。因为那凶残的海盗正是在你刚刚出港时袭击了你，而不是在你满载货物返航之后。正如有人试图制止一头猛狮，却只是擦伤了它的皮肤，并未伤害它，反而更激起它对自己反击，使它日后更加自信和难以捕捉；同样，众人的公敌试图给你重重一击，却击空了，因此使他的对手将来更加警醒和谨慎。\par

2. 人的本性是容易滑倒的，容易受骗，但也容易从欺骗中恢复，正如它迅速跌倒，也容易起来。因为连那位有福之人，我指的是达味，那位被选的君王和先知，在完成了许多善事之后，也承认自己是一个常人；因为他曾迷恋一个外邦女子，而且不止于此，因情欲而犯奸淫，又因奸淫而杀人；但他没有因为已经受了如此沉重的两次打击而试图再给自己第三次打击，而是立即奔向医师，并采用了补救措施：禁食、流泪、哀叹、不断祈祷、频繁告解罪过；这样，他藉这些方式平息了天主的义怒，以致恢复了他先前的地位，以致在奸淫和杀人之后，父亲的功德足以遮盖儿子的偶像崇拜。因为这位达味的儿子，名叫撒罗满，也陷入了与他父亲相同的圈套，因顺从妇女而背离了他先祖的天主。\BibleRef{列上 11:3}你看，不能克制快乐、不能推翻本性中的主导原则、让男人成为女人的奴隶是多么大的邪恶。就是这个撒罗满，他曾是正义而智慧的，但因他的罪而冒失去整个王国的危险，天主因他父亲的声望而允许他保留六分之一的统治权。\BibleRef{列上 11:12}\par

现在，如果你的热忱曾用于世俗的辩论术，然后你因绝望而放弃了它，我会提醒你法庭、审判席以及在那里取得的胜利，还有你从前发言的胆量，并会劝你回到那方面的劳作；但由于我们的竞赛是为天上的事，我们不计较地上的事，我就提醒你另一个法庭，那个可怕而威严的审判席；\BibleQuote{因为我们众人都必须出现在基督的审判台前。}\BibleRef{格后 5:10}那时，那现在被你轻视的一位将坐在审判官的位置上。那时我们该说什么呢？或者我们能提出什么辩护呢，如果我们继续轻视祂？那时我们该说什么呢？我们要以事业的忧虑为由吗？但祂已预先说：\BibleQuote{人纵然赚得了全世界，却赔上了自己的灵魂，为他有什么益处？}\BibleRef{玛 16:26}或者说我们被他人欺骗了？但亚当在辩护时，以妻子为托词并没有帮助他，他说：\BibleQuote{你赐给我的女人，她欺骗了我；}\BibleRef{创 3:12}正如蛇对女人也非借口。钟爱的弟铎，那审判台是可怕的，它不需要原告，也不等待证人；因为\BibleQuote{万物在祂面前都是赤裸敞开的}\BibleRef{希 4:13}审判我们的那一位，我们必须不仅为行为，也为思想交账；因为那位审判者能辨明心思和意念。\BibleRef{希 4:12}但也许你会以本性的软弱为借口，说不能承受那轭。这算什么辩护呢？你说你没有力量承受那容易的轭，你不能承担那轻的担子？难道从疲惫中恢复是痛苦和压人的事吗？因为基督正是在召唤我们，说：\BibleQuote{凡劳苦和负重担的，你们都到我这里来，我要使你们安息。你们背起我的轭，跟我学习，因为我是良善心谦的；因为我的轭是容易的，我的担子是轻的。}\BibleRef{玛 11:28}试问，有什么比摆脱忧虑、事业、恐惧、劳苦，站在人生汹涌波涛之外，安居在平静的港湾更轻省的呢？\par

3. 世间万物中，哪一样在你看来最可羡慕、最值得追求？无疑你会说是统治、财富和声誉。然而，与基督徒的自由相比，这些东西是多么可怜啊！因为统治者受制于民众的愤怒、群众的非理性冲动、对更高统治者的恐惧、以及对被统治者的忧虑；昨日的统治者今天就成了平民；因为现世生命与舞台毫无二致，台上一个人扮演国王，另一个扮演将军，又一个扮演士兵，但夜晚来临时，国王不再是国王，统治者不再是统治者，将军不再是将军；同样，在那一天，每个人将按他的行为得应得的赏报，而不是根据他所扮演的外在角色。那么，荣耀是像草上的花一样凋谢的宝贵事物吗？或者财富，其拥有者被宣告为不幸？\BibleQuote{有祸了}我们读到，\BibleQuote{富有的人}，\BibleRef{路 6:24}并且又说：\Quote{祸哉，那些倚靠自己的力量，以众多财富自夸的人！}但基督徒从不会在做王之后变为平民，富裕之后变为贫穷，被尊敬之后变为卑贱；而是即使贫穷也常富有，谦卑自下时反而受高举；没有人能剥夺他的统治权，除非是那些属黑暗权势的统治者中的一个。\par

\Quote{婚姻是正当的}，你说；我也同意这一点。因为经上记载：\BibleQuote{婚姻应受尊重，床榻也应纯洁；但淫乱和奸淫的人，天主必要审判。}\BibleRef{希 13:4}然而你已不可能再遵守婚姻的正当条件。因为那已归属于天上新郎的人，若离弃祂而投靠一位妻子，那行为便是奸淫，即使你千百次称之为婚姻；更甚的是，这比奸淫更严重，比例正如天主大于人。不要让任何人欺骗你，说：\Quote{天主并没有禁止结婚。}我和你一样知道这一点；祂没有禁止结婚，但祂禁止了奸淫，愿你永不陷入婚姻的束缚！为何你惊奇，当天主被轻视时，婚姻竟被视同奸淫？杀戮带来了义，而仁慈却成了比杀戮更受谴责的原因；因为前者合乎天主的意愿，而后者是被禁止的。\Person{丕乃哈斯}因刺杀那行淫的女子及与她同行的奸夫，此事被算作他的义\BibleRef{户 25:7}；但天主的圣人\Person{撒慕尔}，虽整夜哭泣哀悼恳求，仍无法将\Person{撒乌耳}从天主对他颁布的谴责中解救出来，因为他违背天主的计划，救下了他本应杀死的异族君王。\EditorialNote{撒上第十五章}如果仁慈因违背天主而比杀戮更招致谴责，那么当婚姻涉及弃绝基督时，它比奸淫更定人罪，又有什么可惊奇的呢？因为，正如我起初所说，若你是一个平民，没有人会因逃避兵役而控告你；但现在你已不再属于自己，而是为如此伟大的君王服役。因为妻子对自己的身体没有主权，丈夫却有权\BibleRef{格前 7:4}，何况在基督内生活的人更不能对自己的身体行使主权。如今被轻视的那一位，将来要作我们的审判者；你要时常思念祂和那烈火之河：经上记载\BibleQuote{有火河在祂面前奔流}\BibleRef{达 7:10}；因为一旦被祂交与火中，便不可能期望刑罚有尽头。此世不体面的快乐与影子和梦境无异；因为罪恶的行为尚未完成，快乐的境况便已熄灭；而为这些快乐所受的刑罚却无止境。甜蜜只持续片刻，痛苦却永恒。\par

告诉我，世上有什么是稳固的？财富？它往往连晚上都撑不到。还是荣耀？请听一位义人所说：\BibleQuote{我的生命比奔跑者更快。}\BibleRef{约 9:25}因为正如奔跑者未及站定便疾驰而过，这荣耀也尚未真正临到我们便已飞逝。没有比灵魂更宝贵的了；甚至那些陷入极端愚蠢的人也没有忽略这一点；因为一位外教诗人有言：\Quote{灵魂无等价物。}我知道你在与恶者的争战中变得非常软弱；我知道你正站在快感的火焰中央；但若你对敌人说：\Quote{我们不为你的快感服务，我们不向你一切邪恶的根源屈膝；}若你举目向上，救主甚至现在就要抖落火焰，焚烧那些把你扔进火里的人，并在火窑中为你送来云彩、露水和飒飒凉风，使火焰不能触及你的心思或良心。但不要用火自焚。因为围攻者的武器和机械常常无力摧毁城市的堡垒，但城内一两个居民的背叛却可轻易将他们出卖给敌人。如今，若你内心没有任何思想背叛你，即使恶者从外面引来无数机械，他也将徒劳无功。\par

4. 你靠着天主的恩宠，有许多伟大的人物同情你的困境，鼓励你战斗，为你的灵魂战栗——天主的圣人\Person{瓦莱里乌斯}、各方面都像他兄弟的\Person{弗洛伦提乌斯}、以基督的智慧而明智的\Person{波尔菲里乌斯}，以及其他人。他们每日哀伤，不住地为你祈祷；只要你愿意稍微从敌人手中抽身，他们早已得到所求的了。那么，这岂不奇怪：当别人甚至现在还不绝望于你的救恩，反而不断祈祷，希望他们能得回他们的肢体时，你自己一旦跌倒，就不愿再起来，一直俯伏着，几乎向敌人高喊：\Quote{杀了我，击打我，不要留情？}\BibleQuote{跌倒的岂不再起来吗？}\BibleRef{耶 8:4}——这是上主的神谕。但你却与此抗衡，反驳它；因为如果跌倒的人绝望了，就等于说跌倒的人不再起来。我恳求你，不要对自己造成如此大的伤害；不要向我们倾注如此多的忧伤。我并不是说在你尚未完成二十岁的现在，就连在完成许多大事、终生在基督内度过之后，在极晚年经历这攻击，那时也不应当绝望，而应想起在十字架上成义的强盗，以及约在第十一小时工作却领受了全天工资的工人。但是，正如那些在生命尽头跌倒的人，如果头脑清醒，不应该放弃希望；另一方面，靠这希望度日，说\Quote{这里暂且享受生活的甜蜜，之后我工作一小会儿，就领受整日工作的工资}，也是不安全的。因为我记得常听你说——当许多人劝你常去学校时：\Quote{但我若在短时间内使生命以恶终，我如何能到那曾说\InnerQuote{不要迟延归向上主，不要一日复一日}的主面前去呢？}\BibleRef{德 5:8}恢复这想法，并畏惧那窃贼；因为基督用这名号称呼我们的离世，因为它悄然临到我们。思虑生活中的种种焦虑，既有我们个人的，也有与他人共同的：对统治者的恐惧、市民的嫉妒、常常威胁甚至危及生命的危险、劳苦、痛苦、奴性的谄媚——这些连热心的奴隶都不可做，我们劳苦的果实在这世界终止，这是最令人痛苦的事。确实，许多人命中注定无法享受他们劳苦所得，在劳苦和危险中耗尽了壮年，正当他们希望得到回报时，却空手离世。因为如果一个人历经许多危险，完成许多征战，尚且难以自信地面见地上的君王，那么如果他一生为另一个主人生活并战斗，又如何能面见天上的君王呢？\par

5. 你若要我说起妻子、儿女和奴仆的家务操心，取一个极穷或极富的妻子都是坏事：前者损害丈夫的财产，后者损害他的权威和独立。有孩子是痛苦的，没有孩子更痛苦；因为没有孩子，婚姻便成徒然；有了孩子，则要忍受痛苦的束缚。孩子生病，就引起不小的恐惧；若夭折，便有无可安慰的悲伤；在成长的每个阶段，都有各种忧虑、许多恐惧和辛劳。对于家奴的恶行，又有什么可说的？那么，\Person{特奥多禄}，这便是生活吗？当一个人的灵魂被如此多地分散，当他必须服侍如此多人，为如此多人而活，却从不为自己而活？但在我们中间，朋友，这些事一件也不会发生，我请你作证。因为在那短暂的时间里，当你愿意从这世界的波涛中抬起头时，你知道你享受了何等大的喜乐和欢欣。因为除那为基督而活的人外，没有自由的人。他超越一切烦恼；若他不愿伤害自己，无人能伤害他；他是坚不可摧的；他不因财富损失而刺痛，因他已学会我们\BibleQuote{没有带什么到世界上来，也不能带什么去}；\BibleRef{弟前 6:7} 他不被野心或荣耀的渴望所困，因他已学会我们的籍贯是在天上；\BibleRef{斐 3:20} 无人能用辱骂烦扰他，或用打击激怒他；对基督徒而言，唯一的灾祸是不顺从天主；但其他一切，如财产损失、流放、生命危险，他根本不视为痛苦。而所有人都恐惧的离开此世到另一个世界——这对他比生命本身更甜美。因为正如一个人攀上悬崖顶端，凝望大海和航行其上的人，他看见有些人被波浪冲洗，有些人撞上暗礁，有些人匆忙向一个方向，另一些人被狂风驱赶如囚徒，许多人已在水中，有些人仅以双手代替船和舵，许多人漂在一根木板或船体碎片上，还有些人浮尸，景象多种多样、灾祸各异；同样，服侍基督的人从生活的喧嚣和汹涌波涛中抽身，坐在安全而崇高的高地上。因为有什么地位比那只有一种挂虑\BibleQuote{他该如何讨天主的喜悦？}\BibleRef{得前 4:1} 更为崇高和安全？你是否看见那些航行在这海上的人的船难，\Person{特奥多禄}？因此，我恳求你，避开深海，避开汹涌波涛，抓住一个无法被捕获的高处。有复活，有审判，有一个可怕的审判座在我们离开此世后等待我们；\BibleQuote{我们都必须站在基督的审判座前。}\BibleRef{罗 14:10} 地狱之火并非徒然恐吓我们，如此伟大的福乐也并非无目的地为我们预备。此世之事是影子，甚至比影子更虚无，充满许多恐惧、许多危险和极度的束缚。因此，不要既剥夺那个世界，又剥夺这个世界，因为如果你愿意，你可以两者兼得。那些在基督内生活的人将获得此世之事，\Person{保禄}教导我们时说：\BibleQuote{但我愿意你们免去痛苦}；\BibleRef{格前 7:28} 又说\BibleQuote{但我说这话是为你们的益处}。\BibleRef{格前 7:35} 你看见吗？即使在此世，关切主的事的人也胜过已婚的人。一个已经离开此世到另一个世界的人不可能悔改；没有一个运动员在退出比赛、观众散去后还能再战。\par

要常思念这些事，并粉碎恶者的利剑，他借以毁灭许多人。这利剑便是绝望，它使那些已经跌倒的人断绝希望。这是仇敌的强有力武器，他唯一能压制被俘者的方法就是用这条锁链捆绑他们；如果我们愿意，我们必能藉天主的恩宠迅速折断它。我知道我超过了书信的适当长度，但请原谅我；我不是自愿处于这种状况，而是被我的爱与悲伤所迫，因此我勉强写了这封信，尽管许多人曾阻止我。许多人曾对我说：\Quote{停止徒劳劳动，在磐石上播种。}但我没有听他们任何人的。因为我自己说：有希望，只要天主愿意，我的信会成就一些事；但如果我们所不希望的事发生了，我们至少会有免于因沉默而自责的好处，并且我们将不逊于海上的水手：他们看见自己行业的伙伴因船被风和浪打碎而漂在一块木板上时，便收下帆，抛下锚，乘小船努力营救那些人，即使他们是陌生人，仅因他们的灾难而相识。但如果那些人拒绝被救，无人会指责那些试图救他们的人导致他们的毁灭。这就是我所提供的；但我们相信，藉天主的恩宠，你也会尽你的本分，我们将再次看见你在基督羊群中占据显赫位置。请应圣人们的祈祷，愿我们很快迎接你回来，亲爱的朋友，在真健康中安康。若你对我们有任何关心，且未完全将我们逐出记忆，请惠赐回信；因为这样做将给我们极大喜乐。\par

\SourceRecord{Translated by W.R.W. Stephens.；Nicene and Post-Nicene Fathers, First Series；Vol. 9.；Edited by Philip Schaff.；Buffalo, NY: Christian Literature Publishing Co.,；1889.；<http://www.newadvent.org/fathers/1903.htm>.}
